% ZUGA-IC DTR PDF format - Sablon zorunluluklari
% Font: Carlito (Calibri uyumlu acik kaynak)
% Boyut: 11pt govde, 14pt baslik
% Satir araligi: 1.15
% Hizalama: iki yana yasli
% Kenar boslugu: 2.5 cm

\usepackage{fontspec}
\setmainfont{Carlito}

\usepackage[a4paper, margin=2.5cm]{geometry}

\usepackage{setspace}
\setstretch{1.15}

\usepackage{ragged2e}
\justifying

\usepackage{titlesec}
\titleformat{\section}{\Large\bfseries}{\thesection}{1em}{}
\titleformat{\subsection}{\large\bfseries}{\thesubsection}{1em}{}
\titleformat{\subsubsection}{\normalsize\bfseries}{\thesubsubsection}{1em}{}

% Sayfa numaralandirma
\usepackage{fancyhdr}
\pagestyle{fancy}
\fancyhf{}
\fancyfoot[C]{\thepage}
\renewcommand{\headrulewidth}{0pt}

% Tablo destegi
\usepackage{longtable}
\usepackage{booktabs}

% URL ve link renkleri
\usepackage[colorlinks=true, urlcolor=blue, linkcolor=black]{hyperref}

% Kod blok destegi
\usepackage{listings}
\usepackage{xcolor}
\lstset{
    basicstyle=\ttfamily\small,
    breaklines=true,
    backgroundcolor=\color{gray!10},
    frame=single,
    framerule=0.4pt
}

% Turkce karakter desteği
\usepackage{polyglossia}
\setdefaultlanguage{turkish}
